% -------------------------------------------------
% VERSION TRACKER — No modificar este bloque
% Versión:    Tesis Humanizada (Dummies)
% Archivo:    Capitulo1Introduccion/informacionCapitulos.tex
% Última mod:  2026-05-06T16:45:54
% Git commit:  cc9c766f @ main
% Audiencia:   divulgativo
% Verificar antes de editar: bash check_tesis_version.sh overleaf_tesis/Capitulo1Introduccion/informacionCapitulos.tex
% -------------------------------------------------


\section{Información sobre los capítulos}

La Tabla \ref{tab:estructura_capitulos} describe brevemente cada capítulo del documento e indica qué objetivo específico aborda. Los capítulos 1 a 7 corresponden a la propuesta aprobada en el anteproyecto; los capítulos 8 a 14 documentan la ejecución del trabajo en el orden acordado con la dirección del trabajo de grado.

\begin{longtable}{|p{4.3cm}|p{8.7cm}|p{2.3cm}|}
\caption{Estructura de capítulos.}\label{tab:estructura_capitulos}\\
\cline{1-3}
\cellcolor[gray]{0.9}\textbf{Capítulo} & \cellcolor[gray]{0.9}\textbf{Descripción} & \cellcolor[gray]{0.9}\textbf{Objetivo específico}\\
\cline{1-3}
\endfirsthead
\cline{1-3}
\cellcolor[gray]{0.9}\textbf{Capítulo} & \cellcolor[gray]{0.9}\textbf{Descripción} & \cellcolor[gray]{0.9}\textbf{Objetivo específico}\\
\cline{1-3}
\endhead
Capítulo 1: Introducción & Problema, objetivos, resultados esperados y metodología. & Ninguno\\ \cline{1-3}
Capítulo 2: Justificación y alcance & Por qué se desarrolla el proyecto y qué incluye y excluye. & Ninguno\\ \cline{1-3}
Capítulo 3: Marco referencial & Conceptos, tecnologías y trabajos previos en los que se apoya el proyecto. & Ninguno\\ \cline{1-3}
Capítulo 4: Metodología & Tipo de investigación, enfoque ágil y fases de trabajo. & Ninguno\\ \cline{1-3}
Capítulos 5 a 7: Cronograma, presupuesto e impacto ambiental & Planeación, costos e impacto del proyecto. & Ninguno\\ \cline{1-3}
Capítulo 8: Levantamiento de requerimientos & Entrevista con la empresa del caso de estudio, problema evidenciado y requerimientos funcionales y no funcionales. & 1\\ \cline{1-3}
Capítulo 9: Diseño & Arquitectura, selección de herramientas, modelo de datos, diseño del asistente y de la interfaz. & 2\\ \cline{1-3}
Capítulo 10: Implementación & Patrones de diseño en el código, problemas enfrentados, infraestructura y despliegue. & 3\\ \cline{1-3}
Capítulo 11: Pruebas & Pruebas unitarias, de integración, de extremo a extremo y con usuario, con sus resultados. & 4\\ \cline{1-3}
Capítulo 12: Metodología de gestión del proyecto & Manejo de tareas en GitHub Projects, criterios de aceptación, \textit{Definition of Ready} y \textit{Definition of Done}. & 3\\ \cline{1-3}
Capítulo 13: Trabajos futuros & Mejoras y líneas de continuación identificadas. & Ninguno\\ \cline{1-3}
Capítulo 14: Conclusiones & Cumplimiento de cada objetivo con su indicador. & 1 a 4\\ \cline{1-3}
Anexos & Diagramas UML, modelo entidad relación, flujo FEFO y manual de usuario. & 2 y 3\\ \cline{1-3}
\end{longtable}
